\section{Reflected perturbations and pressure comparison}\label{sec:diagonal-positivity}

\subsection{The finite-volume fold inequality}\label{subsec:slab-finite-fold}

We give the local form of the comparison that will be used to obtain
diagonal positivity. Fix diagonal cells of even height width $h$,
and let $O$ be a real trigonometric polynomial in finitely many
curvature values at plaquettes contained in the base cell $Q_0$.
Thus $O$ is bounded and depends only on this open cell. Let $n\ge8$
be divisible by $4h$, and abbreviate $D_n=D(n/2,n,n)$. Its cell
count, and that of the cubic torus, is
\[
 N_n=\frac{n^3}{|Q|}=\frac{2n^3}{h^3}.
\]
The translations between cells form the integer lattice
\[
 \Lambda_h=\frac h2\Z e_-+\frac h2\Z e_++h\Z e_3.
\]
On $C(n)$ form the translated copies, and on $D_n$ the reflected
copies:
\begin{equation}\label{eq:slab-two-arrays}
 \widetilde O_n(F)=\sum_{v\in\Lambda_h/n\Z^3}O(\tau_v^*F),
 \qquad O_n^{\mathrm{fold}}(F)=\sum_{\kappa\bmod\Gamma_{D_n}}
                                                O(g_\kappa^*F).
\end{equation}
Each sum has $N_n$ terms. The first sum is read on cubic-periodic
curvatures, the second on $D_n$-periodic curvatures. Here $g_\kappa$
is the reflected map from the base cell to $Q_\kappa$.

Let $\mu_{C(n)}$ and $\mu_{D_n}$ be the normalized curvature laws.
For $t>0$ define the tilted diagonal law
\begin{equation}\label{eq:slab-tilted-law}
 \mu_{D_n}^t(\Phi)=
 \frac{\mu_{D_n}(e^{tO_n^{\mathrm{fold}}}\Phi)}
      {\mu_{D_n}(e^{tO_n^{\mathrm{fold}}})}.
\end{equation}
The folded array is invariant under the group generated by the cell
walls: that group permutes the cells of $D_n$ and their reflected
copies. Thus the tilt retains these symmetries.

\begin{proposition}\label{prop:slab-finite-fold}
For the preceding data and every $t>0$,
\begin{equation}\label{eq:slab-finite-fold}
 \mu_{C(n)}(O)\le\mu_{D_n}^t(O)+\frac{|Q|\Delta(n)}{t}.
\end{equation}
The error is the fixed-parameter pressure error of
\eqref{eq:slab-pressure-error}.
\end{proposition}

\begin{proof}
Jensen's inequality and translation invariance on $C(n)$ give
\begin{equation}\label{eq:slab-fold-lower-tangent}
 \log\frac{Z_{C(n)}[e^{t\widetilde O_n}]}{Z_{C(n)}[1]}
                     \ge t\mu_{C(n)}(\widetilde O_n)
                     =tN_n\mu_{C(n)}(O).
\end{equation}
The diagonal chessboard inequality gives
\begin{equation}\label{eq:slab-fold-change}
 Z_{C(n)}[e^{t\widetilde O_n}]
               \le\|e^{tO}\|_{D_n}^{N_n}
               =Z_{D_n}[e^{tO_n^{\mathrm{fold}}}].
\end{equation}
To check that its $N_n$ factors have the same seminorm, translate
the base observable by $v\in\Lambda_h$. Translation permutes the cells
and conjugates their reflected maps, so the fold of the translated
observable is the translate of the original full fold. The law on
$D_n$ is translation invariant, hence its partition function and
the seminorm are unchanged. All observables here are bounded scalar
functions supported in distinct open cells, so the chessboard
inequality applies directly.

Set $L(t)=\log\mu_{D_n}(e^{tO_n^{\mathrm{fold}}})$. Boundedness
permits differentiation, giving
$L'(t)=\mu_{D_n}^t(O_n^{\mathrm{fold}})$. Convexity and $L(0)=0$
give $L(t)\le tL'(t)$. Wall-group invariance of the tilted law makes
the expectation of every reflected copy equal to that of $O$; hence
\begin{equation}\label{eq:slab-fold-upper-tangent}
 \log\frac{Z_{D_n}[e^{tO_n^{\mathrm{fold}}}]}{Z_{D_n}[1]}
                          \le tN_n\mu_{D_n}^t(O).
\end{equation}
Combining \eqref{eq:slab-fold-lower-tangent}--\eqref{eq:slab-fold-upper-tangent}
leaves, in addition, the normalization difference
\[
 \log Z_{D_n}[1]-\log Z_{C(n)}[1]\le n^3\Delta(n).
\]
Division by $tN_n$ proves \eqref{eq:slab-finite-fold}.
\end{proof}

For fixed cell geometry and $t>0$, the pressure error tends to
zero along admissible sides. To remove the tilt we need uniform local
bounds and enough translation symmetry to place an arbitrary observable
next to a prescribed diagonal wall. We prove the bounds in
\cref{sec:local-amplification}. The layer reduction in
\cref{sec:coordinate-symmetry} then recovers coordinate link reflections
and hence even translations in each comparison state. These are the
symmetries used in the final convex separation.

Fix $q>0$, $0<\rho<1$, positive weights $w_i$ and a real mass $M$.
All volume limits in this section keep these parameters fixed. For the comparison construction we
use sides $m_i\to\infty$ such that eight divides all sufficiently
large sides and the set
\[
 \mathcal D=\{m\in\N:m\mid m_i\text{ for all sufficiently large }i\}
\]
is infinite. The set $\mathcal D$ is closed under divisors and least
common multiples. The comparison sequence may be chosen so that every positive integer
divides all sufficiently large sides; \cref{lem:companions-exist}
constructs such sides below any prescribed original sequence.

\subsection{The tilted law at its walls}\label{subsec:diag-wall-law}

Fix diagonal cells with offsets $(c_0,c_+,c_3)$, even width $h$, and
$c_0,c_+\in\Z$, $c_3\in\Z+\tfrac12$, and $4h\in\mathcal D$.
Let $O$ be the real trigonometric polynomial used in the finite fold
comparison. Its support is contained in a box of radius $r_O$, and
we require $h>8r_O+8$. Such widths can be chosen because
$\mathcal D$ contains arbitrarily large multiples of eight. For
large divisible sides, put $D_i=D(m_i/2,m_i,m_i)$ and
$\mu_i^t=\mu_{D_i}^t$, as in \eqref{eq:slab-tilted-law}. Write
\[
 O(F)=\sum_\beta a_\beta e^{i\langle\beta,F\rangle},
 \qquad |O|_1=\sum_\beta|a_\beta|.
\]
The sum is over the distinct exponents in the chosen finite expansion.
In particular $|O|\le|O|_1$. Let $V_i=O_{m_i}^{\mathrm{fold}}$ and
$N_i=m_i^3/|Q|$. For an admissible Grassmann integrand we use
the same notation for the normalized functional
$\mu_i^t(G)=Z_{D_i}[e^{tV_i}G]/Z_{D_i}[e^{tV_i}]$.
On bosonic integrands it agrees with integration against the
tilted probability law. The normalizing denominator satisfies
\begin{equation}\label{eq:diag-tilt-denominator}
 Z_{D_i}[e^{tV_i}]\ge e^{-tN_i|O|_1}Z_{D_i}[1]>0,
 \qquad t\ge0.
\end{equation}
This uses only the positive bosonic measure; it makes no positivity
assertion about a Grassmann integrand.

At a wall $\Pi=(\ell^a,o_a+jh)$, the adapted torus has height period
$p_a$, with $(p_1,p_2,p_3)=(m_i,2m_i,m_i)$. The number $p_a/h$ is
even. Split the cells into the two half-slabs adjacent to this wall,
and define
\[
 \Phi_+(F)=\exp\left(t\sum_{Q_\kappa\subset
              \{o_a+jh<\ell^a<o_a+jh+p_a/2\}}O(g_\kappa^*F)\right),
\]
with representatives in one half-period. Reflection in the wall
maps these cells onto the other half-period modulo the normal closing
vector. Because the reflected copies of $O$ form one wall-invariant
array, on periodic curvatures
\begin{equation}\label{eq:diag-wall-factorization}
 e^{tV_i}=\Phi_+\Theta_\Pi\Phi_+.
\end{equation}
The function $\Phi_+$ is bounded, real, and supported strictly within
the upper half-slab. The support lies there because each plaquette
of a copy of $O$ is contained in its open cell.

Apply the adapted reflected representation to $\Phi_+G$ and
$\Phi_+G'$. For local even integrands $G,G'$ in the upper half-slab,
\begin{equation}\label{eq:diag-tilted-positive-form}
 \mu_i^t(\Theta_\Pi G'\,G)=
 \frac{\operatorname{Tr}
          (\mathfrak S[\Phi_+G']^*\mathfrak S[\Phi_+G])}
      {Z_{D_i}[e^{tV_i}]}.
\end{equation}
The numerator is the Hilbert--Schmidt inner product of the two slab
operators. Thus, on $D_i$, the tilted functional is reflection positive
at each cell wall for local even integrands supported in the corresponding
open positive half-period slab. Wall reflection also gives
$\mu_i^t(\Theta_\Pi G)=\overline{\mu_i^t(G)}$: the unperturbed
functional is reflection covariant, the scalar weight $e^{tV_i}$
is invariant, and its denominator is real and positive.

The bosonic law is invariant under the full group generated by the
cell walls. Indeed, that group preserves $D_i$ and permutes the
reflected copies in $V_i$, while the unperturbed bosonic law is
invariant under its lattice isometries. In particular the tilt is
invariant under translations by $he_-$, $he_+$ and $2he_3$.
These are the periods of the reflected cell array, rather than all
lattice translations.

We use the Polish closed-cochain space $\Omega$ with its product
topology. The pullback by a lattice reflection or
translation is continuous, so invariance passes to weak limits.
Positivity passes for each bounded continuous local positive-side test
once its support lies in the finite positive half-period slab for all
large volumes. Positivity against bounded Borel local tests follows once reflection invariance
has passed: on their finite-dimensional marginal approximate the
test in $L^2$ by compactly supported continuous functions. The
reflected-form difference is at most
\[
 \|\Phi-\Phi_k\|_2(\|\Phi\|_2+\|\Phi_k\|_2),
\]
because reflection preserves the $L^2$ norm. The Gaussian domination
bound also passes to weak limits, by truncating
$e^{\langle\beta,F\rangle}$ and then using monotone convergence.
The conditional-expectation approximation in
\cref{subsec:coordinate-limits} then extends positivity to the whole
open-side $L^2$ space. These facts will be applied after obtaining tightness of the tilted
laws; their extensive weight alone gives no uniform moment estimate.

